% ---- Related Work -----------------------------------------------------------
% Target length: ~1.5 pages. Organise into subsections by approach family.
% Cite, describe briefly, and state how this work differs or builds on it.

\subsection{Foundational Vector Graphics Generation}

Early work on learned vector graphics generation treated SVG as a sequence of
drawing commands. \citet{lopes2019svgvae} introduced SVG-VAE, a class-conditioned
variational autoencoder operating on sequential LSTM representations of Bézier
curve commands, trained on 14 million font characters. While limited to simple,
single-glyph shapes, SVG-VAE established the command-sequence paradigm that
most subsequent work adopts.

\citet{carlier2020deepsvg} extended this with DeepSVG, a hierarchical Transformer
that disentangles high-level shape identity from low-level path commands. DeepSVG
introduced a large-scale SVG icon dataset and demonstrated smooth interpolation
and animation of icon sequences, establishing the icon domain as a productive
testbed for vector generation research.

\subsection{Diffusion-Based Text-to-SVG}

The dominant paradigm from 2023--2024 applies Score Distillation Sampling (SDS)
to optimize Bézier path parameters guided by a pretrained text-to-image diffusion
model. VectorFusion \citep{jain2023vectorfusion} demonstrated that a differentiable
rasterizer (DiffVG) can bridge the pixel and vector domains, allowing gradients
from Stable Diffusion to directly update SVG control points. SVGDreamer
\citep{xing2024svgdreamer} improved upon this with Variational Score Distillation
and semantic-driven path initialization, addressing the over-smoothing and
color saturation issues endemic to SDS-based approaches.

These methods produce visually plausible results for illustrations and logos but
are poorly suited to icon generation for three reasons: (1) generation takes
30--120 seconds per image; (2) the optimized paths do not respect icon design
conventions; and (3) the resulting SVG is typically an over-parameterized
approximation rather than a compact, editable structure.

\subsection{LLM-Native SVG Code Generation}

A parallel line of work treats SVG generation as code generation and leverages
the sequence modeling capabilities of large language models. This is the paradigm
our work follows.

IconShop \citep{wu2023iconshop} tokenizes SVG path data into a discrete vocabulary
and trains an autoregressive Transformer on a curated icon corpus. It supports
text-guided generation, interpolation between icons, and semantic combination.
While closely related to our approach, IconShop uses a purpose-built tokenization
scheme and is not built on a pretrained code language model.

StarVector \citep{rodriguez2025starvector} treats image-to-SVG vectorization as
code generation using a Vision-Language Model backbone. Trained on the 2M-sample
SVG-Stack corpus, it handles diverse SVG types including icons, logos, and
technical diagrams. LLM4SVG \citep{xing2025llm4svg} introduces 55 domain-specific
SVG semantic tokens and trains on 250K SVGs with 580K instruction pairs,
demonstrating that structured tokenization improves generation quality.

OmniSVG \citep{yang2025omnisvg} unifies image-conditioned and text-conditioned
generation via a Qwen2.5-VL backbone and the 2M-sample MMSVG-2M dataset. SVGen
\citep{svgen2025} applies Chain-of-Thought reasoning and Group Relative Policy
Optimization (GRPO) reinforcement learning, showing that explicit reasoning
steps significantly improve structural validity.

\paragraph{Distinction from our work.}
All the above LLM-based methods are trained on general SVG corpora and
evaluated on diverse SVG types. We focus exclusively on the icon domain, which
admits a narrower distribution (single-color fills, 1--20 paths, 16--512px
viewBox) and therefore benefits from a specialized training set. We build the
first ML training corpus from the Iconify ecosystem and perform
domain-specific fine-tuning on Qwen3.5-9B, targeting the quality and
structural conventions expected of production icons.

\subsection{Automated SVG Captioning and Datasets}

Several recent works construct large SVG training datasets using automated
annotation pipelines. StarVector \citep{rodriguez2025starvector} renders SVG
files to images and uses a VLM to generate text descriptions, assembling the
SVG-Stack corpus. LLM4SVG \citep{xing2025llm4svg} similarly uses automated
pipelines to produce 580K SVG-text instruction pairs from 250K SVGs. SVGen
\citep{svgen2025} applied this to approximately 500K icons from the Iconfont
repository to produce the SVG-1M dataset.

Our \dataset{} corpus applies the same automated annotation strategy to the
Iconify ecosystem. The key differences are: (1) a rigorous per-set license
filter preserving \num{275912} icons under permissive open-source terms;
(2) use of structured short-and-long caption pairs tuned for instruction
following rather than single descriptions; (3) full documentation of the
normalization pipeline to ensure reproducibility.

\subsection{Evaluation of SVG Generation}

SVGenius \citep{tang2025svgenius} benchmarks 22 LLMs on SVG understanding,
editing, and generation tasks, finding that reasoning-enhanced models outperform
pure scaling and that quality degrades sharply with rising SVG complexity.
VectorGym \citep{vectorgym2025} provides human-annotated benchmarks across four
SVG tasks. Both works confirm the need for domain-specific evaluation metrics
beyond standard image quality scores. Our evaluation framework adapts these
insights to the icon domain.
